%% Movement arrows, \mvfrom and \mvto, on the page.
%%
%% An arrow is nothing but vector ink, so everything here is measured on a
%% rendering rather than read off the text layer:
%%
%%   MVPLAIN  one arrow below the line, from BASEA to LANDA.  The HEAD is at
%%            the landing site -- the direction is the whole meaning of the
%%            arrow, and an arrow drawn backwards sits at the right
%%            coordinates.  And the line after it is pushed down to make
%%            room: MVCTRL is the same pair of lines with no arrow.
%%   KLAND    the room goes PAST the arrow, not just down to it: in a list
%%            of sub-examples the next item follows with no skip between,
%%            and a line set only \lineskip under the horizontal reads as
%%            struck through.  HLAND is the same above the line, under an
%%            item in capitals, which have no depth to keep it off.
%%   MVNEST   an arrow over a stretch that holds another goes a level
%%            further out, so the two horizontals do not run into one.
%%   CHMID    a word that is both a landing site and a base position (a
%%            chain): the arrow arriving and the one leaving meet it side by
%%            side, not down one line.
%%   GLAND    in a gloss the arrow goes ABOVE the object line, since below
%%            it would cut through the gloss tier.
%%   ULAND    [above] outside a gloss.
%%   LONELY   an \mvto with no partner draws nothing and says so -- and
%%            the \mvfrom{z} in the NEXT example is not its partner.
%%   DOG      the ends start at one depth below a word with a descender
%%            and one without: the strut is the floor, not the ink.
%%   MVUP     the same above the line: an arrow over "ace" starts above
%%            where the ascender of "fly" ends.
%%   MLAND    level=3 puts the horizontal two steps further out.
%%   SLAND    an arrow above and one below never push each other out a
%%            level, however their stretches overlap.
%%   WLAND    two ends on two lines: the arrow still joins them, and the
%%            log says the room was only kept on one line.
%%   RUNL     a mark that starts a paragraph starts it, rather than being
%%            set as a line of its own.
%%   PLAND    two ends on two pages: nothing is drawn, and the log says so.
%%
%% Two passes: the positions come from the previous run's .aux, and the
%% room is reserved in the run itself, so the second run is final -- which
%% the absence of the rerun warning pins.
\documentclass[11pt]{article}
\usepackage[a4paper,margin=25mm]{geometry}
\usepackage{iftex}
\ifPDFTeX
  \usepackage[T1]{fontenc}
  \usepackage[utf8]{inputenc}
\else
  \usepackage{fontspec}
\fi
\usepackage{linguexx}
\pagestyle{empty}


\begin{document}
\ex. MVPLAIN \mvto{a}{LANDA} stays here \mvfrom{a}{BASEA} end.

\ex. MVNEXT the line under an arrow.

\ex. MVCTRL no arrow here.

\ex. MVCTRLNEXT the line under no arrow.

\ex. \a. MVSUB \mvto{k}{KLAND} stays here \mvfrom{k}{KBASE} end.
     \b. KNEXT HANGS UNDER THE ARROW.
     \z.

\ex. \a. KPREV SITS OVER THE ARROW.
     \b. MVSUBUP \mvto[above]{h}{HLAND} stays here \mvfrom{h}{HBASE} end.
     \z.

\ex. MVNEST \mvto{o}{OUTL} \mvto{i}{INL} x \mvfrom{i}{INB} y \mvfrom{o}{OUTB} z.

\ex. \mvto{b}{CHW} do you think \mvto{a}{\mvfrom{b}{CHMID}} John bought
     \mvfrom{a}{CHEND}?

\ex. MVPREV the line above a gloss.

\ex. \gll \mvto{g}{GLAND} hat Hans \mvfrom{g}{GBASE} gekauft\\
          what has Hans GTIER bought\\

\ex. MVABOVE \mvto[above]{u}{ULAND} x \mvfrom{u}{UBASE}.

\ex. MVLONE \mvto{z}{LONELY} has no partner.

\ex. MVLONETWO and \mvfrom{z}{LONETWO} is not it: a label is the example's own.

\ex. MVDESC \mvto{d}{DOG} x \mvfrom{d}{gypsy}.

\ex. MVLEVEL \mvto[level=3]{m}{MLAND} x \mvfrom{m}{MBASE}.

\ex. MVSIDES \mvto{s}{SLAND} \mvto[above]{t}{TLAND} x \mvfrom{t}{TBASE} y
     \mvfrom{s}{SBASE}.

\ex. MVUP \mvto[above]{e}{ace} x \mvfrom{e}{fly} end.

\ex. MVWRAP a line that ends in \mvto{w}{WLAND}\linebreak
     \mvfrom{w}{WBASE} short.

\mvto{r}{RUNL} starts a paragraph of running text \mvfrom{r}{RUNB} here.

\ex. MVPAGE \mvto{p}{PLAND} before the break.\par
\newpage
     PAFTER after it, \mvfrom{p}{PBASE} on the next page.
\end{document}
